% ==================== 2 物理来源、选择定则与互斥规则 ====================
\section{物理来源、选择定则与互斥规则}

\subsection{物理来源}

两者都源于分子振动能级跃迁，但物理过程不同。红外属吸收过程：振动必须改变分子偶极矩（$\Delta\mu\neq0$），光子能量正好等于振动能级差时被吸收。拉曼属散射过程：激光把分子激发到虚能级后几乎立即返回，若返回能级与出发时不同，散射光便产生波数差（拉曼位移）；强度取决于极化率变化（$\Delta\alpha\neq0$），常温下斯托克斯线远强于反斯托克斯线（图~\ref{fig:energy}）。

\begin{figure}[H]
  \centering
  \begin{tikzpicture}[font=\footnotesize,>=Latex,line width=0.6pt]
    % (a) 红外吸收
    \node at (1.2,4.15) {(a) 红外吸收};
    \draw (0,0) -- (2.4,0);
    \draw (0,1.6) -- (2.4,1.6);
    \draw[dashed] (0,3.2) -- (2.4,3.2);
    \node[left] at (0,0) {$v=0$};
    \node[left] at (0,1.6) {$v=1$};
    \node[left] at (0,3.2) {虚能级};
    \draw[->,thick] (1.2,0) -- (1.2,1.6);
    \node[right] at (1.25,0.8) {$h\nu_0=\Delta E$};
    % (b) 拉曼散射
    \node at (7.5,4.15) {(b) 拉曼散射};
    \draw (4.5,0) -- (10.5,0);
    \draw (4.5,1.6) -- (10.5,1.6);
    \draw[dashed] (4.5,3.2) -- (10.5,3.2);
    \node[left] at (4.5,0) {$v=0$};
    \node[left] at (4.5,1.6) {$v=1$};
    \node[left] at (4.5,3.2) {虚能级};
    % 瑞利
    \draw[->] (5.4,0) -- (5.4,3.2);
    \draw[->] (5.6,3.2) -- (5.6,0);
    \node[below] at (5.5,-0.05) {瑞利};
    % 斯托克斯
    \draw[->] (7.0,0) -- (7.0,3.2);
    \draw[->] (7.2,3.2) -- (7.2,1.6);
    \node[below] at (7.1,-0.05) {斯托克斯};
    \node[right] at (7.25,2.4) {$h(\nu_0-\nu_{\mathrm{vib}})$};
    % 反斯托克斯
    \draw[->] (8.9,1.6) -- (8.9,3.2);
    \draw[->] (9.1,3.2) -- (9.1,0);
    \node[below] at (9.0,-0.05) {反斯托克斯};
    \node[right] at (9.15,2.4) {$h(\nu_0+\nu_{\mathrm{vib}})$};
    \node[right] at (5.65,2.2) {$h\nu_0$};
  \end{tikzpicture}
  \caption{红外吸收与拉曼散射的能级过程。}
  \label{fig:energy}
\end{figure}

\subsection{选择定则与互斥规则的适用条件}
\noindent\textbf{1. 选择定则。}群论中，IR 活性要求振动模式的不可约表示与偶极矩分量（$x$、$y$、$z$）的不可约表示之一相同；Raman 活性则要求该不可约表示与极化率张量分量（$x^2$、$y^2$、$z^2$、$xy$、$yz$、$zx$）的不可约表示之一相同。

\noindent\textbf{2. 互斥规则。}它只在有对称中心时成立：中心对称群的振动只能属于 g（对反演对称）或 u（反对称）两类，偶极矩属 u 型、极化率属 g 型，同一振动不能同时具有两种活性。

\noindent\textbf{3. 适用条件。}结构确有对称中心，如 CO$_2$（$D_{\infty h}$）、锐钛矿（$I4_1/amd$）与金红石（$P4_2/mnm$）；比较对象为同一振动模式；结构为理想周期，表面、缺陷与吸附会降低局域对称性，使禁戒模弱激活\cite{taudul2023}。

CO$_2$ 对称伸缩 1333 cm$^{-1}$ 只见于拉曼、反对称伸缩 2349 cm$^{-1}$ 只见于红外\cite{larkin2017}；TiO$_2$ 的 E$_g$ 144/447 cm$^{-1}$ 属 g 模，u 模落在 400 cm$^{-1}$ 以下远红外。油酸无对称中心，两谱均可观察，但极性 COO$^-$ 的偶极矩变化大、极化率变化小，红外强度远高于拉曼\cite{ning2014}。所以红外主要判断配位，拉曼主要判断晶格与晶相。
